\section{Fortitude: Cardinal Virtue and Its Parts}
\label{sec:fortitude}

Fortitude perfects the irascible appetite where a rational good is hard to attain or defend. Its primary act is endurance: attack is valuable when needed, but steadfastness under grave evil most directly prevents fear from driving reason away from the good. Aristotle treats courage and several virtues of aspiration and expenditure in \emph{Nicomachean Ethics} III--IV; Aquinas orders fortitude and its annexed virtues in ST II--II, qq.~123--140.

\cardinalgroup{The principal virtue: endurance and attack}

\virtuedossier
  {Fortitude (courage)}
  {Principal cardinal virtue; Aristotle, \emph{Nicomachean Ethics} II.7 and III.6--9; Augustine, \emph{On the Morals of the Catholic Church} I.15.25; ST II--II, qq.~123--127; CCC 1808.}
  {Fear and daring before grave dangers, paradigmatically death faced for a just and proportionately great good.}
  {To stand firm and, when reason requires, to attack danger rather than abandon the good through fear.}
  {\textbf{Cowardice} yields to fear against right reason (\oppositioncode{M}).}
  {\textbf{Rashness} rushes into danger through disordered daring or deficient fear (\oppositioncode{M}).}
  {Fearlessness from ignorance, anger, compulsion, experience, or confidence in superior force can resemble courage without being its act. Martyrdom is a chief act of fortitude, not a separate moral virtue.}

Fortitude neither suppresses all fear nor seeks danger as proof of identity. Fear can accurately register a great evil; virtue keeps that perception within the hierarchy of goods. A prudent retreat, protective concealment, or appeal for help can be courageous when it preserves the true end. A spectacular exposure that serves vanity or needlessly endangers dependents is rash even when the agent feels no fear.

\cardinalgroup{Quasi-integral requirements in fortitude's proper matter}

\begin{studybox}{Confidence and security}
In the most difficult matter of fortitude, \textbf{confidence} supports attack and \textbf{security} steadies the person against fear. Aquinas calls them quasi-integral in fortitude's own principal matter while reducing their secondary matters under magnanimity and related virtues (ST II--II, q.~128, a.~1). They are requirements or movements of complete fortitude here, not two additional coordinate habits.
\end{studybox}

\cardinalgroup{Annexed virtues for undertaking difficult goods}

\virtuedossier
  {Magnanimity}
  {Annexed virtue of fortitude and Aristotelian mean; Aristotle, \emph{Nicomachean Ethics} II.7 and IV.3; ST II--II, qq.~129--133.}
  {Great honor as testimony to great moral excellence and to difficult works proportioned to one's God-given powers.}
  {To undertake what is truly great and worthy while judging one's capacity truthfully and ordering honor to virtue.}
  {\textbf{Pusillanimity} shrinks from great things within one's capacity (\oppositioncode{M}).}
  {Three source-named excesses differ by respect: \textbf{presumption concerning capacity} undertakes beyond one's powers; \textbf{ambition} desires honor inordinately; and \textbf{vainglory} desires glory inordinately (\oppositioncode{M}).}
  {One rational mean can be exceeded under several circumstances or formal respects. Magnanimity is compatible with humility: humility restrains claims beyond truth before God; magnanimity refuses to bury gifts entrusted for great service. It is not dominance, celebrity, or self-sufficiency.}

\virtuedossier
  {Magnificence}
  {Annexed virtue of fortitude and Aristotelian mean; Aristotle, \emph{Nicomachean Ethics} II.7 and IV.2; ST II--II, qq.~134--135.}
  {Large expenditure ordered to a proportionately great, durable, or public work.}
  {To conceive and complete a great work with expenditure fitting its end, materials, circumstances, and agent.}
  {\textbf{Parvificence} (Latin \emph{parvificentia}, rendered paltriness or meanness) spoils a great work by spending beneath its fitting scale or attending to petty savings (\oppositioncode{M}).}
  {\textbf{Wasteful consumption} is Aquinas's excess; Aristotle's \emph{banausia}/\emph{apeirokalia} is often rendered vulgar or tasteless display (\oppositioncode{M}). Both exceed proportion to work, person, and end.}
  {Magnificence differs from liberality: both use wealth, but magnificence concerns the greatness of a work rather than ordinary giving. Cost alone cannot make a project magnificent.}

\cardinalgroup{Annexed virtues for enduring difficulty}

\virtuedossier
  {Patience}
  {Annexed virtue of fortitude; Augustine, \emph{On Patience}; ST II--II, q.~136.}
  {Sorrows and evils endured while remaining faithful to a rational and, in infused virtue, supernatural good.}
  {To prevent sorrow from overthrowing reason or causing abandonment of the good.}
  {\textbf{Impatience} is the named contrary state (\oppositioncode{K}), but Aquinas does not construct a dedicated two-flank mean for patience.}
  {No intrinsic excess-side vice exists (\oppositioncode{NA}). Apathy toward another's suffering is a defect of perception or charity, not excessive patience.}
  {Patience is not passivity, approval of abuse, or refusal of proportionate remedy. One can oppose an evil patiently by refusing hatred and remaining governed by the good.}

Augustine distinguishes genuine patience by its end: enduring hardship for greed, pride, or crime displays toughness but not virtue. Christian patience receives suffering without treating suffering itself as the supreme good. It can lament, seek justice, protect the vulnerable, and persevere without surrendering reason to sorrow.

\virtuedossier
  {Longanimity}
  {Thomistic virtue formally allied most closely with magnanimity, yet comprised under patience insofar as delay causes sorrow; ST II--II, q.~136, a.~5.}
  {A worthy good whose attainment is distant in time and therefore wearisome to desire.}
  {To sustain the movement of hope toward a delayed good without abandoning the work that prepares for it.}
  {\textbf{Impatience at delay} is expressly contrary, but Aquinas establishes no dedicated vice species unique to longanimity (\oppositioncode{U}).}
  {No intrinsic excess-side vice exists (\oppositioncode{NA}); procrastination chooses or causes delay through neglect, whereas longanimity bears a delay.}
  {Because it tends through hope toward a distant good, longanimity has more in common with magnanimity than patience. It is nevertheless comprised under patience insofar as delay produces sorrow, and it differs from perseverance's endurance of prolonged difficulty in action.}

\virtuedossier
  {Perseverance}
  {Annexed virtue of fortitude; Aristotle, \emph{Nicomachean Ethics} VII.7; ST II--II, qq.~137--138.}
  {The duration and accumulated difficulty of continuing in a good work to its fitting completion.}
  {To remain in the good course for as long as reason and the end require.}
  {\textbf{Softness} (Latin \emph{mollities}) abandons the good because one cannot bear ordinary toil or loss of pleasure (\oppositioncode{M}).}
  {\textbf{Pertinacity} persists when reason requires stopping, changing means, or yielding a private position (\oppositioncode{M}).}
  {The historical technical sense of \emph{mollities} concerns surrender to difficulty and carries no judgment about sex, disability, or gentle temperament. Perseverance is not stubborn attachment to the original tactic.}

\virtuedossier
  {Constancy}
  {Subordinate specification that pertains to perseverance; ST II--II, q.~137, a.~3.}
  {External obstacles and changing pressures that threaten a settled good purpose.}
  {To keep practical commitment stable when opposition, fortune, or surrounding conditions push it away from reason.}
  {Abandoning a purpose under external pressure opposes constancy, but Aquinas supplies no dedicated vice noun for this specification (\oppositioncode{U}); prudential inconstancy has a different formal relation.}
  {\textbf{Obstinacy} counterfeits stability by refusing correction when reason changes the judgment (\oppositioncode{X}); it is not excess constancy.}
  {Perseverance primarily answers long duration; constancy answers external impediments. Neither requires retaining an erroneous plan after relevant facts or duties change.}

\cardinalgroup{Acts and semblances that are not additional virtues}

Confidence, preparedness, aggression, and endurance are integral movements within fortitude's field, not automatically separate virtues. Confidence without realistic self-knowledge becomes presumption; endurance without a good end becomes hardness; attack without authority or proportion becomes violence. Magnanimity and magnificence extend the family's reach from resisting evil to undertaking difficult positive goods. Patience, longanimity, perseverance, and constancy distinguish why remaining faithful is hard: present sorrow, delayed attainment, prolonged toil, or external obstruction.
